\Chapter{Opis problemu i analiza obecnej sytuacji}\label{chapter:problem_analysis}

\section{Kontekst projektu}
\subsection{Cyfryzacja ochrony zdrowia w Unii Europejskiej}
W ostatnich latach cyfryzacja stała się jednym z głównych kierunków rozwoju systemów ochrony zdrowia w krajach Unii Europejskiej. Rozwiązania z zakresu eHealth są stopniowo wdrażane zarówno w dużych szpitalach, jak i mniejszych placówkach medycznych, a elektroniczna dokumentacja medyczna zaczyna pełnić rolę podstawowego źródła informacji o pacjencie. Proces ten jest częścią szerszej strategii transformacji cyfrowej UE, której celem jest zwiększenie dostępności usług publicznych oraz poprawa jakości funkcjonowania instytucji ochrony zdrowia \cite{EU_eHealth_2024}.

Jednym z najważniejszych założeń programu „Digital Decade 2030” jest zapewnienie wszystkim obywatelom Unii Europejskiej dostępu do elektronicznej dokumentacji medycznej do roku 2030 \cite{EU_eHealth_2024}. W związku z tym państwa członkowskie rozwijają centralne systemy umożliwiające przechowywanie i udostępnianie danych zdrowotnych w formie cyfrowej. Według raportu „2024 Digital Decade eHealth Indicator Study” średni poziom rozwoju usług eHealth w krajach UE osiągnął w 2023 roku wartość 79\%, co oznacza dalszy wzrost względem poprzednich lat \cite{EU_eHealth_2024}. Polska została wskazana jako jedno z państw o wysokim poziomie dojrzałości cyfrowych usług zdrowotnych, osiągając wynik 90\% \cite{EU_eHealth_2024}.

Raport pokazuje również, że większość państw Unii Europejskiej wykorzystuje scentralizowane platformy dostępu do danych medycznych, a coraz większa liczba usług dostępna jest zarówno przez portale internetowe, jak i aplikacje mobilne \cite{EU_eHealth_2024}. Oznacza to, że współczesne systemy ochrony zdrowia odchodzą od modelu opartego wyłącznie na dokumentacji papierowej i stopniowo przechodzą do środowiska cyfrowego, w którym informacje mogą być dostępne w czasie rzeczywistym dla uprawnionych użytkowników.

Rozwój systemów eHealth wynika również z potrzeby usprawnienia przepływu informacji pomiędzy różnymi jednostkami ochrony zdrowia. W raportach UE podkreśla się, że cyfrowe systemy medyczne mają poprawiać dostęp do danych, wspierać ciągłość leczenia oraz zwiększać zaangażowanie pacjentów w proces opieki zdrowotnej \cite{EU_eHealth_2024}. Szczególnie istotne staje się zapewnienie możliwości szybkiego dostępu do wyników badań, historii leczenia oraz danych dotyczących terapii i przyjmowanych leków.

Pomimo dynamicznego rozwoju technologii wiele krajów nadal mierzy się z problemami związanymi z interoperacyjnością systemów, integracją danych oraz dostępnością niektórych typów dokumentacji medycznej \cite{EU_eHealth_2024}. Raport wskazuje między innymi, że część danych, takich jak obrazy medyczne czy informacje dotyczące implantów, nadal nie jest powszechnie dostępna dla obywateli w formie elektronicznej \cite{EU_eHealth_2024}. Istotnym wyzwaniem pozostaje również zapewnienie odpowiedniego poziomu bezpieczeństwa danych, zgodności z regulacjami UE oraz dostępu wyłącznie dla uprawnionych osób.

Cyfryzacja ochrony zdrowia wiąże się także z koniecznością uwzględnienia potrzeb różnych grup społecznych. W dokumentach UE podkreśla się znaczenie tworzenia systemów dostępnych dla osób starszych, osób z ograniczonymi kompetencjami cyfrowymi oraz użytkowników wymagających dodatkowego wsparcia technicznego \cite{EU_eHealth_2024}. Z tego powodu nowoczesne systemy medyczne powinny być projektowane nie tylko jako zaawansowane technologicznie rozwiązania, ale również jako narzędzia intuicyjne, czytelne i możliwie proste w codziennym użytkowaniu.

\subsection{Cel i założenia projektu}

Projekt dotyczy stworzenia zintegrowanego systemu wspierającego zarządzanie informacjami o pacjentach w Szpitalu im. Bajgla we Wrocławiu. Placówka planuje wdrożenie rozwiązania umożliwiającego gromadzenie, aktualizowanie oraz udostępnianie dokumentacji medycznej w ramach jednego spójnego środowiska cyfrowego. Głównym celem projektu jest usprawnienie przepływu informacji pomiędzy oddziałami oraz ograniczenie problemów wynikających z funkcjonowania dokumentacji papierowej.

Założeniem projektu jest stworzenie systemu umożliwiającego personelowi medycznemu szybki dostęp do najważniejszych informacji o pacjencie, historii leczenia, wyników badań oraz aktualnych zleceń medycznych. Rozwiązanie ma wspierać codzienną pracę lekarzy, pielęgniarek, laboratoriów i administracji poprzez uporządkowanie obiegu danych oraz ograniczenie konieczności ręcznego przekazywania dokumentacji pomiędzy oddziałami.

Projekt uwzględnia również potrzeby pacjentów związane z łatwiejszym dostępem do własnej dokumentacji medycznej oraz możliwością sprawdzania wyników badań w formie elektronicznej. System ma poprawić komunikację pomiędzy placówką a pacjentem oraz ograniczyć konieczność korzystania z dokumentacji papierowej.

Jednym z kluczowych założeń projektu jest także zapewnienie wysokiego poziomu bezpieczeństwa danych medycznych. System powinien umożliwiać kontrolę dostępu do informacji, rozdzielenie uprawnień użytkowników oraz ochronę danych zgodnie z obowiązującymi regulacjami dotyczącymi prywatności i bezpieczeństwa informacji w Unii Europejskiej. Ważnym aspektem projektu pozostaje również stworzenie rozwiązania intuicyjnego i prostego w codziennym użytkowaniu, tak aby ograniczyć problemy związane z wdrażaniem nowych technologii w środowisku szpitalnym.
\section{Analiza obecnej sytuacji}
Szpital im. Bajgla we Wrocławiu jest dużą, wielooddziałową placówką medyczną. W strukturze szpitala znajdują się między innymi oddziały chirurgiczne, interna, neurologia, ortopedia, onkologia, pediatria, oddział intensywnej terapii, Szpitalny Oddział Ratunkowy, laboratorium diagnostyczne, pracownia diagnostyki obrazowej oraz prosektorium. Ze względu na skalę działalności placówka każdego dnia obsługuje dużą liczbę pacjentów, co wiąże się z koniecznością ciągłego przetwarzania i przekazywania ogromnej ilości informacji medycznych pomiędzy różnymi działami szpitala.

Obecnie dokumentacja medyczna w Szpitalu im. Bajgla prowadzona jest głównie w formie papierowej. Karty pacjentów, wyniki badań, opisy diagnostyczne, zlecenia lekarskie, raporty pielęgniarskie oraz dokumenty administracyjne przechowywane są fizycznie i przekazywane pomiędzy oddziałami ręcznie przez personel medyczny lub administracyjny. Powoduje to sytuacje, w których dostęp do informacji jest ograniczony czasowo i zależy od aktualnego miejsca przechowywania dokumentacji. W praktyce oznacza to, że personel często musi samodzielnie odnajdywać potrzebne dokumenty, kontaktować się z innymi działami lub oczekiwać na dostarczenie wyników badań i raportów.

Szczególnie problematyczny jest przepływ informacji pomiędzy oddziałami a laboratorium i diagnostyką obrazową. Wyniki badań laboratoryjnych oraz opisy badań obrazowych przekazywane są w sposób tradycyjny, co wydłuża czas obiegu dokumentacji i utrudnia bieżący dostęp do aktualnych danych pacjenta. Stanowi to krytyczny problem w sytuacjach wymagających szybkiej reakcji.

Brak centralnego systemu cyfrowego powoduje również problemy związane z archiwizacją dokumentacji. Dane pacjentów przechowywane są w wielu segregatorach, kartotekach i archiwach papierowych, przez co odnalezienie starszej historii leczenia może wymagać dużej ilości czasu. W przypadku pacjentów hospitalizowanych wielokrotnie personel musi ręcznie przeszukiwać dokumentację z poprzednich pobytów, aby uzyskać pełny obraz przebiegu leczenia.

Obecny sposób funkcjonowania placówki wpływa bezpośrednio na codzienną pracę lekarzy, pielęgniarek oraz pracowników administracyjnych. Znaczna część czasu poświęcana jest nie na kontakt z pacjentem, lecz na organizację dokumentacji, przekazywanie danych oraz uzupełnianie formularzy papierowych. W wielu przypadkach te same informacje przepisywane są kilkukrotnie do różnych dokumentów, co zwiększa ryzyko błędów i nieścisłości. Problem nasila się szczególnie podczas dużego obciążenia oddziałów lub nagłych sytuacji wymagających szybkiego dostępu do informacji o stanie zdrowia pacjenta.

Chaos informacyjny widoczny jest także z perspektywy pacjentów. Osoby hospitalizowane nie mają bezpośredniego dostępu do swoich wyników badań ani informacji o przebiegu leczenia. Uzyskanie dokumentacji wymaga kontaktu z personelem lub odbioru dokumentów w formie papierowej.

Dodatkowym wyzwaniem jest rosnąca liczba danych medycznych generowanych każdego dnia przez szpital. Rozwój diagnostyki, większa liczba wykonywanych badań oraz rosnąca liczba pacjentów powodują, że tradycyjny model przechowywania dokumentacji staje się coraz mniej efektywny organizacyjnie. Obecny sposób pracy nie odpowiada również kierunkowi rozwoju nowoczesnej ochrony zdrowia w Unii Europejskiej, w której coraz większy nacisk kładzie się na cyfryzację usług medycznych, szybki dostęp do danych oraz integrację informacji o pacjencie \cite{EU_eHealth_2024}.

Analiza obecnej sytuacji pokazuje, że głównym problemem Szpitala im. Bajgla jest brak jednego, centralnego środowiska umożliwiającego sprawne zarządzanie dokumentacją medyczną i przepływem informacji pomiędzy użytkownikami. 
\subsection{Perspektywa dyrektora i organizacji szpitala}

Dyrekcja Szpitala im. Bajgla dostrzega, że obecny sposób organizacji dokumentacji medycznej przestaje być wystarczający dla placówki o tak dużej skali działalności. Wraz ze wzrostem liczby pacjentów oraz ilości wykonywanych badań coraz trudniejsze staje się skuteczne zarządzanie przepływem informacji. Z perspektywy dyrektora szczególnie istotne staje się usprawnienie organizacji pracy i ograniczenie sytuacji, w których personel medyczny poświęca znaczną część czasu na czynności administracyjne.

Istotnym problemem dla organizacji szpitala jest także brak jednolitego standardu pracy z dokumentacją medyczną. Poszczególne oddziały wypracowały własne sposoby przekazywania informacji i organizowania dokumentów, przez co przepływ danych nie jest spójny w całej placówce.

Dyrekcja szpitala zwraca również uwagę na rosnące wymagania związane z dalszą cyfryzacją ochrony zdrowia w Unii Europejskiej. Placówka musi stopniowo dostosowywać się do nowych standardów dotyczących elektronicznej dokumentacji medycznej, bezpieczeństwa danych oraz dostępności usług cyfrowych dla pacjentów \cite{EU_eHealth_2024}.

Dyrekcja oczekuje rozwiązania, które pozwoli personelowi szybciej uzyskiwać dostęp do danych pacjentów, monitorować przebieg leczenia oraz ograniczyć problemy wynikające z opóźnionego przekazywania dokumentacji. Szczególne znaczenie ma to w przypadku oddziałów o dużej dynamice pracy, takich jak SOR czy oddział intensywnej terapii, gdzie szybki dostęp do informacji może mieć bezpośredni wpływ na bezpieczeństwo pacjentów.

Jednocześnie dyrekcja jest świadoma, że wdrożenie nowego systemu będzie dużą zmianą organizacyjną dla całego szpitala. Istnieje obawa, że nieodpowiednio zaprojektowane rozwiązanie mogłoby zwiększyć liczbę obowiązków administracyjnych lub spotkać się z oporem części pracowników. 
\subsection{Perspektywa personelu medycznego}
Personel medyczny Szpitala im. Bajgla należy do grupy użytkowników, na których obecny sposób organizacji dokumentacji wpływa najbardziej bezpośrednio. Lekarze, pielęgniarki, ratownicy medyczni, technicy laboratoryjni oraz pracownicy diagnostyki każdego dnia uczestniczą w procesie przekazywania dużej ilości informacji dotyczących pacjentów. 

Z perspektywy lekarzy jednym z największych problemów jest brak szybkiego dostępu do pełnej historii leczenia pacjenta. Uzyskanie wszystkich potrzebnych informacji wymaga często przeglądania wielu dokumentów, kontaktowania się z innymi oddziałami lub oczekiwania na dostarczenie wyników badań. Lekarze zwracają uwagę, że czas poświęcany na organizację dokumentacji ogranicza możliwość skupienia się bezpośrednio na pracy z pacjentem.

Podobne problemy dotyczą personelu pielęgniarskiego, który oprócz obowiązków związanych z opieką nad pacjentami odpowiada również za prowadzenie i aktualizowanie części dokumentacji medycznej. Powoduje to dodatkowe obciążenie organizacyjne oraz zwiększa ryzyko pomyłek wynikających z nieczytelnych zapisów, zagubionych dokumentów lub nieaktualnych danych.

Istotnym problemem dla personelu jest również przepływ informacji dotyczących badań laboratoryjnych i diagnostycznych. Pracownicy laboratoriów oraz diagnostyki obrazowej wykonują dużą liczbę badań, jednak przekazywanie wyników w tradycyjnej formie wydłuża cały proces komunikacji. Personel medyczny nie posiada możliwości bieżącego monitorowania statusu badań, przez co konieczne staje się wielokrotne sprawdzanie dostępności wyników lub kontaktowanie się z innymi działami szpitala. 

Część pracowników obawia się cyfryzacji placówki. Wśród najczęściej pojawiających się problemów znajdują się obawy przed koniecznością nauki nowych systemów, zmiany dotychczasowych przyzwyczajeń oraz zwiększenia liczby obowiązków administracyjnych. Niektórzy pracownicy postrzegają wdrożenie nowego rozwiązania jako potencjalne utrudnienie codziennej pracy, szczególnie jeśli system byłby skomplikowany. Widoczny jest również strach przed popełnianiem błędów podczas korzystania z nowych narzędzi cyfrowych.

Jednocześnie część personelu dostrzega korzyści wynikające z wprowadzenia centralnego systemu informatycznego. Lekarze i pielęgniarki wskazują przede wszystkim potrzebę szybkiego dostępu do dokumentacji, łatwiejszego przekazywania informacji oraz ograniczenia ilości pracy papierowej. Personel oczekuje rozwiązania, które pozwoli uprościć codzienne obowiązki, zmniejszyć liczbę powtarzalnych czynności administracyjnych oraz usprawnić komunikację pomiędzy oddziałami.

Z punktu widzenia użytkowników bardzo ważna staje się intuicyjność systemu. Nowe rozwiązanie powinno być możliwie proste i czytelne. System nie może wymagać skomplikowanej obsługi ani dużej liczby dodatkowych działań podczas wykonywania podstawowych obowiązków. Istotne znaczenie ma również mobilność rozwiązania, szczególnie w przypadku pracy na oddziałach oraz podczas przemieszczania się pomiędzy pacjentami. Możliwość szybkiego dostępu do danych z poziomu urządzeń mobilnych mogłaby znacząco usprawnić organizację pracy personelu.

Pracownicy oczekują rozwiązania, które rzeczywiście uprości wykonywanie obowiązków, ograniczy chaos organizacyjny oraz poprawi komunikację pomiędzy oddziałami, a nie kolejnego systemu zwiększającego poziom skomplikowania pracy.
\subsection{Perspektywa pacjenta}
Pacjenci Szpitala im. Bajgla są grupą bezpośrednio odczuwającą skutki problemów związanych z organizacją dokumentacji w placówce. Obecny model funkcjonowania szpitala sprawia, że pacjent ma ograniczony dostęp do danych dotyczących własnego leczenia oraz przebiegu hospitalizacji. Uzyskanie wyników badań, historii leczenia czy informacji o dalszych etapach terapii wymaga najczęściej kontaktu z personelem lub odbioru dokumentów w formie papierowej.
Szczególnie problematyczne staje się to w przypadku pacjentów korzystających z usług kilku oddziałów jednocześnie lub wymagających długotrwałego leczenia. Informacje dotyczące badań, konsultacji oraz zaleceń lekarskich są rozproszone, przez co pacjent często nie posiada pełnego obrazu własnego procesu leczenia. Z perspektywy pacjentów ważnym problemem jest czas oczekiwania na informacje. Odbiór wyników badań lub dokumentacji wymaga często osobistej wizyty w szpitalu albo długiego oczekiwania na przekazanie danych przez personel. W praktyce utrudnia to kontakt pacjenta z własną dokumentacją medyczną oraz ogranicza możliwość szybkiego reagowania na zmiany związane ze stanem zdrowia.

Pacjenci zwracają również uwagę na brak przejrzystości procesu leczenia. Osoby hospitalizowane nie zawsze wiedzą, na jakim etapie znajdują się wykonane badania, kiedy pojawią się wyniki oraz jakie działania będą realizowane w dalszej kolejności.
Jednocześnie pacjenci coraz częściej oczekują rozwiązań cyfrowych podobnych do tych wykorzystywanych w innych usługach publicznych i prywatnych. Możliwość sprawdzenia wyników badań online, łatwiejszy kontakt z placówką oraz szybszy dostęp do dokumentacji medycznej postrzegane są jako istotne udogodnienia poprawiające komfort korzystania z opieki zdrowotnej. W raportach dotyczących rozwoju eHealth podkreśla się, że zwiększenie dostępu obywateli do własnych danych medycznych stanowi jeden z głównych kierunków rozwoju nowoczesnej ochrony zdrowia \cite{EU_eHealth_2024}.
\subsection{Bezpieczeństwo danych i dostęp}
Jednym z najważniejszych wyzwań związanych z funkcjonowaniem Szpitala im. Bajgla jest zapewnienie odpowiedniego poziomu bezpieczeństwa dokumentacji medycznej oraz kontrola dostępu do informacji o pacjentach. Dane zdrowotne należą do kategorii danych szczególnie wrażliwych, dlatego sposób ich przechowywania, przetwarzania oraz udostępniania wymaga zachowania wysokich standardów bezpieczeństwa i zgodności z obowiązującymi regulacjami prawnymi.

W obecnym modelu dostęp do danych zależy od fizycznej lokalizacji dokumentów. Taki model utrudnia skuteczną kontrolę nad tym, kto w danym momencie korzysta z dokumentacji oraz czy wszystkie informacje są odpowiednio zabezpieczone przed nieuprawnionym dostępem. Problem staje się szczególnie widoczny w dużej placówce, gdzie każdego dnia z dokumentacją pracuje wiele osób.

Istotnym zagrożeniem jest również możliwość zagubienia dokumentów lub ich nieprawidłowego przekazania pomiędzy oddziałami. Dokumentacja papierowa może zostać przypadkowo pozostawiona w niewłaściwym miejscu, przekazana nieodpowiedniej osobie lub ulec uszkodzeniu.

Z perspektywy organizacji szpitala ważne staje się również rozdzielenie poziomów dostępu do danych. Poszczególni pracownicy nie potrzebują dostępu do wszystkich informacji znajdujących się w dokumentacji medycznej. Lekarze, pielęgniarki, laboranci, administracja czy pracownicy prosektorium korzystają z różnych zakresów danych związanych z wykonywanymi obowiązkami. W obecnym modelu papierowym kontrola takich uprawnień jest utrudniona, ponieważ dokumentacja może być fizycznie przekazywana pomiędzy wieloma osobami bez możliwości dokładnego monitorowania dostępu.

Ważnym aspektem są również wymagania wynikające z regulacji obowiązujących na terenie Unii Europejskiej. Cyfrowe systemy ochrony zdrowia muszą być zgodne między innymi z zasadami GDPR dotyczącymi ochrony danych osobowych oraz bezpieczeństwa informacji \cite{EU_eHealth_2024}. Oznacza to konieczność zapewnienia odpowiedniego uwierzytelniania użytkowników, kontroli dostępu do dokumentacji oraz możliwości monitorowania działań wykonywanych w systemie. 
Istotnym wyzwaniem pozostaje również zapewnienie bezpieczeństwa przy jednoczesnym zachowaniu prostoty użytkowania systemu. Personel medyczny pracuje pod dużą presją czasu, dlatego procedury związane z logowaniem i dostępem do danych nie mogą nadmiernie utrudniać codziennej pracy.
\subsection{Problemy techniczne i organizacyjne}
Planowana cyfryzacja Szpitala im. Bajgla wiąże się nie tylko z koniecznością stworzenia nowego systemu informatycznego, ale również z dużą zmianą organizacyjną obejmującą sposób funkcjonowania całej placówki. Ze względu na skalę szpitala oraz liczbę oddziałów wdrożenie nowego rozwiązania wymaga uwzględnienia wielu problemów technicznych, logistycznych i organizacyjnych, które mogą wpływać na późniejsze działanie systemu oraz jego akceptację przez użytkowników.

Jednym z głównych wyzwań jest konieczność zaprojektowania systemu dostosowanego do pracy wielu różnych działów szpitala jednocześnie. Każdy oddział funkcjonuje w nieco inny sposób i posiada własne potrzeby związane z dokumentacją oraz przepływem informacji. Oznacza to, że system musi być wystarczająco elastyczny, aby umożliwić obsługę różnych procesów medycznych bez nadmiernego komplikowania interfejsu użytkownika.

Istotnym problemem technicznym jest również stworzenie odpowiedniej infrastruktury umożliwiającej funkcjonowanie systemu w całym szpitalu. Placówka wymaga stabilnego dostępu do sieci, odpowiednich urządzeń dla personelu medycznego oraz bezpiecznego środowiska przechowywania danych. 
Wdrożenie systemu obejmującego wiele oddziałów oznacza konieczność zapewnienia ciągłości działania nawet w przypadku awarii części infrastruktury lub problemów technicznych. 

Ważnym wyzwaniem jest także przeniesienie istniejącej dokumentacji papierowej do nowego środowiska cyfrowego. Szpital posiada dużą ilość archiwalnych danych dotyczących pacjentów, które przez lata były przechowywane w formie tradycyjnej. Proces cyfryzacji dokumentacji wymaga więc ustalenia sposobu migracji danych, określenia priorytetów oraz przygotowania procedur umożliwiających stopniowe przechodzenie na nowy model pracy bez paraliżowania funkcjonowania placówki.

Problemem pozostaje również przygotowanie personelu do pracy z nowym systemem. Wdrożenie systemu nie może ograniczać się wyłącznie do aspektów technicznych. Konieczne jest przeprowadzenie szkoleń, przygotowanie materiałów pomocniczych oraz stworzenie procedur wsparcia dla użytkowników na początkowym etapie działania systemu.

Dodatkowym wyzwaniem organizacyjnym jest konieczność zachowania ciągłości pracy szpitala podczas wdrażania nowego rozwiązania. Placówka funkcjonuje całodobowo i nie może zostać czasowo wyłączona z działania na potrzeby instalacji systemu. W praktyce oznacza to konieczność etapowego wdrażania poszczególnych funkcjonalności oraz równoczesnego funkcjonowania starego i nowego sposobu pracy przez pewien okres przejściowy.

Ważnym aspektem pozostaje także skalowalność rozwiązania. System powinien umożliwiać dalszy rozwój placówki, obsługę rosnącej liczby pacjentów oraz możliwość rozbudowy o nowe funkcjonalności w przyszłości.
\section{Podsumowanie analizy problemu}
\subsection{Podsumowanie problemu}
Analiza funkcjonowania Szpitala im. Bajgla pokazuje, że głównym problemem placówki jest brak centralnego systemu umożliwiającego sprawne zarządzanie dokumentacją medyczną oraz przepływem informacji pomiędzy oddziałami. Obecny model pracy oparty głównie na dokumentacji papierowej powoduje utrudniony dostęp do danych pacjentów, wydłuża proces przekazywania informacji oraz zwiększa obciążenie administracyjne personelu medycznego. Problemy te wpływają zarówno na organizację pracy szpitala, jak i komfort oraz bezpieczeństwo pacjentów.

Istotnym problemem jest brak jednolitego środowiska pracy dla wszystkich działów placówki. Obecny sposób organizacji pracy nie odpowiada także rosnącym wymaganiom związanym z cyfryzacją ochrony zdrowia oraz potrzebą szybkiego dostępu do danych medycznych w czasie rzeczywistym.

Projektowane rozwiązanie musi uwzględniać potrzeby różnych grup użytkowników. System powinien wspierać codzienną pracę lekarzy, pielęgniarek, laboratoriów oraz administracji, jednocześnie pozostając rozwiązaniem prostym i intuicyjnym w obsłudze. 

Ważnym aspektem projektu pozostaje także zapewnienie odpowiedniego poziomu bezpieczeństwa danych medycznych oraz stworzenie przejrzystego modelu dostępu do informacji w zależności od roli użytkownika.

\subsection{Docelowa funkcjonalność}
Za sukces projektu zostanie uznane stworzenie systemu umożliwiającego cyfrową obsługę dokumentacji pacjentów w Szpitalu im. Bajgla bez konieczności wykorzystywania papierowego obiegu danych w codziennej pracy personelu medycznego. System musi umożliwiać dodawanie, przeglądanie oraz aktualizowanie dokumentacji medycznej w czasie rzeczywistym przez uprawnionych użytkowników z różnych oddziałów szpitala.

Kryterium sukcesu będzie również zapewnienie centralnego dostępu do danych pacjentów z poziomu jednego środowiska informatycznego. Lekarz po zalogowaniu do systemu powinien mieć możliwość sprawdzenia historii leczenia pacjenta, wyników badań laboratoryjnych, dokumentacji diagnostycznej oraz aktualnych zleceń bez konieczności korzystania z dodatkowych źródeł informacji.

System musi umożliwiać obsługę różnych poziomów dostępu do danych. Każda grupa użytkowników powinna posiadać wyłącznie uprawnienia odpowiadające wykonywanym obowiązkom. Oznacza to między innymi, że laboratorium może dodawać wyniki badań, lekarz posiada dostęp do pełnej dokumentacji medycznej pacjenta, natomiast administracja nie ma możliwości przeglądania danych medycznych niewymaganych do realizacji obowiązków organizacyjnych.

Ważnym kryterium sukcesu będzie również możliwość dostępu pacjentów do wybranych informacji medycznych online. Pacjent po zalogowaniu powinien mieć możliwość sprawdzenia wyników badań, historii wizyt oraz podstawowych informacji dotyczących przebiegu leczenia bez konieczności odbierania dokumentacji w formie papierowej.
Projekt zostanie uznany za skuteczny, jeśli system będzie umożliwiał jednoczesną pracę wielu użytkowników bez utraty danych oraz zapewni ciągłość działania w środowisku całodobowej pracy szpitala. 